\emph{(i)} With $\nInf\cdot\nO = 1$ and $Q(\nInf) = Q(\nO) = 0$, the vector $X(s) = \nO + sn$ is null for every $s$, since $n$ is null and orthogonal to $\nInf$ and $\nO$; it is the conformal point of $x = sn$ (the general point $\nO + x - \tfrac12 Q(x)\nInf$ with $Q(x) = 0$). The rotor $T_{t} = 1 + \tfrac{t}{2}n\nInf$ satisfies $T_{t}\tilde T_{t} = 1$ because $(n\nInf)^{2} = -Q(n)Q(\nInf) = 0$; likewise $K_{c} = 1 + \tfrac{c}{2}m\nO$, and $D_{\theta} = \cosh\theta + \sinh\theta\,\nu e_{-}$ with $(\nu e_{-})^{2} = 1$. The sandwiches $T_{t}X(s)\tilde T_{t}$, $K_{c}X(s)\tilde K_{c}$, $D_{\theta}X(s)\tilde D_{\theta}$ are computed exactly in the check at rational $s, t, c$ and $e^{\theta}$ and are nonzero multiples of $X(s + t)$, $X(s/(1 - 2cs))$ (here $m\cdot n = 2$) and $X(e^{-2\theta}s)$. These are the Moebius maps of the matrices $\begin{pmatrix}1 & t\\ 0 & 1\end{pmatrix}$, $\begin{pmatrix}1 & 0\\ -2c & 1\end{pmatrix}$, $\begin{pmatrix}e^{-\theta} & 0\\ 0 & e^{\theta}\end{pmatrix}$, which generate $SL(2, \R)$; the product $\begin{pmatrix}1 & -1\\ 0 & 1\end{pmatrix}\begin{pmatrix}1 & 0\\ 1 & 1\end{pmatrix}\begin{pmatrix}1 & -1\\ 0 & 1\end{pmatrix} = \begin{pmatrix}0 & -1\\ 1 & 0\end{pmatrix}$ is realised by $J = T_{-1}K_{-1/2}T_{-1}$, and the check confirms $JX(s)\tilde J \propto X(-1/s)$ and $J\nO\tilde J \propto n$. Since $X(s) \propto s^{-1}\nO + n \to n$ projectively, $n$ is the point at infinity of the line, and $J$ exchanges it with $s = 0$: the null line is a circle on which the rotors of the attachment (Theorem \claimref{synthesis.attachment-conformal-energy-is-generated-by-z}) act by Moebius maps. This realisation is on null lines of Minkowski space; what the rest of the proof uses is only the $SL(2, \R)$ covariance of the vacuum kernel below, which holds on any generator (on the eternal hole the point at infinity is not a spacetime point joined by the attachment). For the kernel, $g(u) - g(v) = (u - v)/((cu + d)(cv + d))$ when $ad - bc = 1$, so $(cu + d)^{-1}(cv + d)^{-1}(g(u) - g(v))^{-1} = (u - v)^{-1}$; the weight $(cu + d)^{-1}$ is $g'(u)^{1/2}$ up to the sign fixed by the spin structure, the weight-$\tfrac12$ action of Theorem \claimref{synthesis.local-normality-at-the-edge-gives-the-null-ground-state}(i), and the map $\psi(u) \mapsto (cu + d)^{-1}\psi(g(u))$ is unitary on $L^{2}(\R)$. So $\omega_{0}$ is invariant, and in the coordinate $w = -1/u$ near $\infty$ its kernel has the same form $\tfrac{I_{4}}{2\pi}(w - w' + I_{4}0)^{-1}$ as near any point. Kernel differences transform by the same unitary, so square integrability near $\infty$ is posed in $w$ exactly as near a finite point, and on $\R$ it is square integrability on the arcs through $\infty$, the complements of bounded intervals.

\emph{(ii)} Let $E_{1} = \{|u| > L\}$ and $E_{2}, E_{3}$ disjoint open bounded intervals inside $[-L, L]$, and $A_{k} = \R\setminus E_{k}$: $A_{1} = [-L, L]$ is a bounded interval and $A_{2}, A_{3}$ are complements of bounded intervals, so all three are intervals of the closed generator. A point lies in at most one $E_{k}$, so a pair $(x, y)$ meets at most two of them and avoids the third, $k$ say, whence $(x, y) \in A_{k}\times A_{k}$. Hence $\R\times\R = \bigcup_{k}A_{k}\times A_{k}$ and $\|R\|^{2}_{L^{2}(\R^{2})} \le \sum_{k}\|R\|^{2}_{L^{2}(A_{k}^{2})} < \infty$. The check confirms the cover on a rational grid with points on both far ends, and that pairs with one point beyond $-L$ and one beyond $L$ lie only in the arcs through $\infty$. Local normality on $A_{k}$ gives the Hilbert-Schmidt condition there by the argument of Theorem \claimref{synthesis.local-normality-at-the-edge-gives-the-null-ground-state}(ii) (Powers-Stormer for the compressed symbols, $\|A - B\|_{\mathrm{HS}} \le 2\|A^{1/2} - B^{1/2}\|_{\mathrm{HS}}$); for $A_{2}, A_{3}$ the algebra is that of $L^{2}(A_{k})\otimes K$ and the argument is unchanged. If $W_{\omega} = P$ is a projection, Fact \claimref{synthesis.imported-powers-stormer-takesaki-disjointness}(a) with $A = P$, $B = P_{0}$ gives quasi-equivalence exactly when $P^{1/2} - P_{0}^{1/2} = P - P_{0}$ and $(1 - P)^{1/2} - (1 - P_{0})^{1/2} = -(P - P_{0})$ are Hilbert-Schmidt, that is, when $R$ is square integrable on $\R^{2}$; pure states that are quasi-equivalent are unitarily equivalent. Conversely, quasi-equivalence of the global states gives quasi-equivalence of their restrictions to the algebra of each $A_{k}$ (restrictions of gauge-invariant quasi-free states are again factor states, and normal states of one factor are quasi-equivalent), hence the condition on every interval.

\emph{(iii)} If $\omega$ is translation-stationary, $R(u, v) = r(u - v)$, and $\int_{L}^{L + N}\int_{L}^{L + N}|r(u - v)|^{2}\,du\,dv \ge N\int_{0}^{1}\int_{0}^{1}|r(u - v)|^{2}\,du\,dv$, the $N$ diagonal unit squares having equal integrals by translation invariance (checked exactly for polynomial $r$ with exact rectangle integrals); square integrability on any far end therefore forces $r = 0$, so $\omega = \omega_{0}$. If $\omega$ is stationary for the boost about a point $u_{0}$ (a corollary of the stored theorem, not a new statement), which acts on the generator as the weight-$\tfrac12$ dilation about $u_{0}$ (Theorem \claimref{synthesis.stationary-horizon-state-is-its-own-edge-scaling-limit}(i)), the condition on an interval containing $u_{0}$ gives $\omega = \omega_{0}$ by Theorem \claimref{synthesis.local-normality-at-the-edge-gives-the-null-ground-state}(iv) (stated about the edge; a translation moves $u_{0}$ there and preserves $\omega_{0}$). In both cases $\omega$ satisfies (N$_{+}$) (Theorem \claimref{synthesis.null-vacuum-on-a-horizon-generator}(ii)).

\emph{(iv)} By (ii) a pure $\omega$ with the condition is unitarily equivalent to $\omega_{0}$, so $\omega = \langle\Psi, \pi_{0}(\cdot)\Psi\rangle$ for a unit vector $\Psi$ of the representation $\pi_{0}$ of $\omega_{0}$. There the affine translations $\tau_{b}$ are implemented by $U(b) = e^{I_{4}bP}$ with $P \ge 0$, $\omega_{0}$ being the ground state (Theorem \claimref{synthesis.null-vacuum-on-a-horizon-generator}(ii)), and $\pi_{0}(\tau_{b}(a)) = U(b)\pi_{0}(a)U(b)^{*}$ is the implementation in the representation of $\omega$ as well. Hence $\omega(P) = \langle\Psi, P\Psi\rangle \ge 0$ for $\Psi$ in the form domain of $P$. For the last sentence: for a unit vector $g \in L^{2}(\R)\otimes K$ in the range of $1 - P_{0}$ (the empty modes of $\omega_{0}$, in the convention where $W_{0}$ is the projection $P_{0}$), the vector $\pi_{0}(\psi^{*}(g))\Omega_{0}$ gives the pure gauge-invariant quasi-free state with projection $P_{0} + |g\rangle\langle g|$; its kernel differs from $W_{0}$ by the rank-one kernel $g(u)\overline{g(v)}$, square integrable and not translation invariant (a translation-invariant square-integrable kernel on $\R^{2}$ is zero, by the estimate of (iii)). So it satisfies the condition, is quasi-equivalent to $\omega_{0}$, and is not a ground state; any condition depending only on the folium admits it.

\emph{(v)} The kernel of the state KMS at temperature $T$ for the translations is $\tfrac{I_{4}}{2\pi}\,\pi T/\sinh(\pi T(s + I_{4}0))$, $s = u - v$, whose singular part at $s = 0$ is that of $W_{0}$. With $x = \pi T s$, $\pi T/\sinh x = s^{-1}\cdot x/\sinh x$, and the exact inversion of $\sinh x/x = \sum_{k} x^{2k}/(2k + 1)!$ gives $x/\sinh x = 1 - \tfrac{x^{2}}{6} + \tfrac{7x^{4}}{360} - \dots$ (checked to order $8$), so the difference from $W_{0}$ is $\tfrac{I_{4}}{2\pi}s^{-1}(x/\sinh x - 1) = -\tfrac{I_{4}}{2\pi}\big(\tfrac{(\pi T)^{2}}{6}s - \tfrac{7(\pi T)^{4}}{360}s^{3} + \dots\big)$, analytic and nonzero. It is translation invariant, so by the estimate of (iii) it is not square integrable on $(L, \infty)^{2}$ for any $L$: the state fails the Hilbert-Schmidt condition on the arc through $\infty$ (and on either far end), while near each finite point its difference is bounded and the condition holds.

\emph{(vi)} The polynomial $f' = 1 - \tfrac34 S((u + 1)/2)$ on $[-1, 1]$ has $f'(-1) = 1$, $f'(1) = \tfrac14$, $f'' = f''' = 0$ at both ends (since $S'$ and $S''$ vanish at $t = 0, 1$), is decreasing ($S' = 30t^{2}(1 - t)^{2} \ge 0$), hence in $[\tfrac14, 1]$, and $\int_{-1}^{1}f' = 2 - \tfrac34 = \tfrac54 = f(1) - f(-1)$; so $f$ is a $C^{3}$ increasing bijection of $\R$, linear outside $[-1, 1]$ (all checked exactly). The unitary $\psi(u) \mapsto f'(u)^{1/2}\psi(f(u))$ maps the projection $P_{0}$ to $W_{f}$, so $\omega_{f}$ is pure. On $u, v \ge 1$, $\sqrt{f'(u)f'(v)}/(f(u) - f(v)) = \tfrac14/((u - v)/4) = (u - v)^{-1}$, and on $u, v \le -1$ trivially, so $W_{f} = W_{0}$ on both far squares, and the restrictions of $\omega_{f}$ and $\omega_{0}$ to the algebras of $(1, \infty)$ and of $(-\infty, -1)$ coincide (a gauge-invariant quasi-free state is fixed by its two-point kernel). On bounded intervals the difference is continuous across the diagonal by Theorem \claimref{synthesis.null-vacuum-singularity-under-reparametrisation-and-dilation}(i), so the local Hilbert-Schmidt condition holds and, by Theorem \claimref{synthesis.local-normality-at-the-edge-gives-the-null-ground-state}(ii), $\omega_{f}$ is locally the null vacuum at every finite point. On the cross block $u \ge 1$, $v \le -1$: $\sqrt{f'(u)f'(v)} = \tfrac12$ and $f(u) - f(v) = u/4 - v$, so $2\pi(W_{f} - W_{0})/I_{4} = 2/(u - 4v) - 1/(u - v) = (u + 2v)/((u - 4v)(u - v))$, homogeneous of degree $-1$. On the rectangle $[5, 6]\times[-2, -1]$ the numerator is at least $1$ and the factors of the denominator at most $14$ and $8$, so the modulus is at least $\tfrac{1}{112}$; the dilates by $2^{k}$ are disjoint and, the kernel being of degree $-1$ and the area scaling by $4^{k}$, each carries the same lower bound $(1/112)^{2}$ for the integral of the squared modulus (checked exactly for $40$ dilates). So $W_{f} - W_{0}$ is not square integrable on $\R^{2}$, and by (ii) $\omega_{f}$ and $\omega_{0}$ are not quasi-equivalent; being factor states they are disjoint (Fact \claimref{synthesis.imported-powers-stormer-takesaki-disjointness}(a)). For slope $1/\mu$ at $+\infty$ (so $\sqrt{f'(u)f'(v)} = \mu^{-1/2}$ and $f(u) - f(v) = u/\mu - v$ on the cross block) the cross-block difference is $\sqrt{\mu}/(u - \mu v) - 1/(u - v)$, of degree $-1$, and two slopes $1/\mu_{1}$, $1/\mu_{2}$ give the degree $-1$ difference $\sqrt{\mu_{1}}/(u - \mu_{1}v) - \sqrt{\mu_{2}}/(u - \mu_{2}v)$, which for $\mu_{1}, \mu_{2} = 4, 9$ (slopes $\tfrac14$, $\tfrac19$) is $-(u + 6v)/((u - 4v)(u - 9v))$, nonzero; the same estimate shows they are mutually disjoint. In $w = -1/u$, $f$ becomes $w \mapsto w$ for small $w > 0$ (from $u \le -1$) and $w \mapsto 4w$ for small $w < 0$ (from $-1/(u/4) = 4w$): a slope jump at $w = 0$, so the cross block is the degree $-1$ obstruction of Theorem \claimref{synthesis.local-normality-at-the-edge-gives-the-null-ground-state}(v) at the point at infinity.

Implementability. For an affine map $\alpha$, $\omega_{f}\circ\alpha$ is the vacuum of the parameter $f\circ\alpha$; transporting by the unitary of $f^{-1}$, which maps $\omega_{f}$ to $\omega_{0}$, it is quasi-equivalent to $\omega_{f}$ exactly when the vacuum of $h = f\circ\alpha\circ f^{-1}$ is quasi-equivalent to $\omega_{0}$, that is (pure states, (ii)), when $W_{h} - W_{0}$ is square integrable on $\R^{2}$. For $\alpha$ the translation by $b$, $h(U) = U + b$ near $-\infty$ and $h(U) = U + b/4$ near $+\infty$; for $\alpha$ the dilation by $\lambda$ about $0$, $h(U) = \lambda U$ near both ends (checked exactly on the far pieces). In either case $h$ is a regular reparametrisation with the same slope $s$ at both ends. Then $W_{h} - W_{0}$ vanishes on each far square ($W_{0}$ is invariant under affine maps); on the cross far block it is $0$ for the dilation and $-d/((U - V)(U - V + d))$, $d = b/4 - b$, for the translation, of degree $-2$, and $\int_{U, W > L}(U + W)^{-4} = 1/(24L^{2})$ is finite; on bounded sets it is bounded (as above); and for $U$ in a bounded set $B$ and $|V|$ large, $h(V) = sV + c_{\pm}$, so $\sqrt{h'(U)s}/(h(U) - sV - c_{\pm}) - 1/(U - V) = O(1/|V|)$ uniformly in $U \in B$, square integrable on $B$ times a far end. So $W_{h} - W_{0}$ is square integrable and every translation and every dilation about $0$ is implementable in the representation of $\omega_{f}$; the boosts about other points are their composites. The slope jump between the two ends is what no affine map of the generator changes and what the arcs through $\infty$ detect.

\emph{Reading.} Statement (ii) uses the Hilbert-Schmidt form of local normality throughout; that the vacuum of a regular parameter is locally normal to $\omega_{0}$ on bounded intervals in the stronger sense of quasi-equivalence would follow from the implementability of smooth circle diffeomorphisms for the free chiral fermion (Segal; cited from memory, not used). The principle of local normality is Fact \claimref{synthesis.imported-haag-local-normality}, applied to the algebras of intervals of the closed generator; the closing itself is (i).
